\begin{textsample}{POS Dim 1 – persona\_grok – Score 15.00 – t418\_grok.txt}  \label{ex:f1_pos_040}
Migrant fishers from Southeast Asia often \textbf{board} industrial fishing vessels with \textbf{hopes} of securing \textbf{better} lives for themselves and their \textbf{families}.
However, these aspirations frequently give \textbf{way} to brutal realities resembling modern slavery, marked by grueling long \textbf{hours}, physical and verbal abuse, and in some tragic cases, even death at sea.
As consumers, we have the power to make a \textbf{difference} by opting for sustainable seafood choices and spreading \textbf{awareness} about these injustices to pressure the industry for change.
Greenpeace has brought these hidden \textbf{stories} to light through two powerful \textbf{films}, shot in Taiwan and featuring real Indonesian workers.
These \textbf{films} are rooted in true accounts uncovered by Greenpeace East Asia and Southeast Asia.
They portray the grim \textbf{experiences} of fishers, including the unfair treatment they endure and their deep personal motivations, such as providing for their \textbf{families}.
More than 30 fishers were \textbf{interviewed} for this \textbf{work}, many expressing \textbf{fear} of losing their \textbf{jobs} if they speak out against the conditions.

% matched lemmas: awareness, board, difference, experience, family, fear, film, good, hope, hour, interview, job, story, way, work
\end{textsample}
